\section{Problem Formulation}
\label{sec:formulation}

We formalize the supervision selection problem for visual tool-use learning.
We first define the setting, then identify the failure mode of standard outcome-based filtering, and finally state our selection criterion.

\subsection{Setup and Notation}

Consider a multimodal question $q$ with ground-truth answer $a^*$ and an associated image $I$.
A \emph{tool-augmented trajectory} $\tau = (r_1, c_1, o_1, r_2, c_2, o_2, \ldots, r_K, \hat{a})$ interleaves reasoning steps $r_k$, tool calls $c_k$, returned observations $o_k$, and a final answer $\hat{a}$.
We write $\mathcal{O}(\tau) = \{o_1, \ldots, o_K\}$ for the set of all tool observations in $\tau$.

A teacher model $\mathcal{M}_T$ generates candidate trajectories; a student model $\mathcal{M}_S$ is trained on the retained subset.
The goal is to select trajectories that teach $\mathcal{M}_S$ \emph{when and why} to invoke visual tools, not merely \emph{how} to imitate tool-call syntax.

\subsection{The Problem of Spurious Tool Supervision}

Standard trajectory-distillation pipelines retain any trajectory whose final answer is correct, i.e., $\hat{a} = a^*$.
We call this \textbf{outcome-based filtering}.
While it removes obviously broken traces, it admits a pervasive failure mode:

\begin{definition}[Spurious Tool Trajectory]
\label{def:spurious}
A trajectory $\tau$ with $\hat{a} = a^*$ is \emph{spurious} if the teacher model can produce the correct answer without any tool observation, i.e.,
\begin{equation}
    P_{\mathcal{M}_T}(\hat{a} = a^* \mid q, I, \varnothing) \geq P_{\mathcal{M}_T}(\hat{a} = a^* \mid q, I, \mathcal{O}(\tau)),
\end{equation}
where $\varnothing$ denotes the absence of all tool observations.
\end{definition}

Spurious trajectories arise because strong teacher models often already possess sufficient visual understanding to answer without tools; the tool calls they produce are incidental rather than necessary.
Training on such trajectories teaches the student to \emph{correlate} tool invocation with correct answers, but does not teach it to \emph{rely} on the acquired evidence—leading to over-calling, hallucinated observations, or failure to transfer tool-use skills to harder problems where evidence is genuinely needed.

\subsection{Learnable Trajectories}

We argue that a trajectory provides useful supervision if and only if it satisfies two jointly necessary conditions:

\begin{definition}[Learnable Trajectory]
\label{def:learnable}
A trajectory $\tau$ is \emph{learnable} if it satisfies:
\begin{enumerate}[leftmargin=*, itemsep=2pt]
    \item \textbf{Outcome Correctness.} The final answer is correct: $\hat{a} = a^*$.
    \item \textbf{Causal Observation Utility.} The tool observations are causally necessary for reaching the correct answer (Definition~\ref{def:causal}).
\end{enumerate}
\end{definition}

Condition~1 is standard; it ensures the trajectory demonstrates a successful reasoning path.
Condition~2 is the key addition: it requires that the observations acquired via tool calls actually \emph{make a difference} to the outcome.

\subsection{Causal Observation Utility}

We operationalize causal necessity through a counterfactual test inspired by interventionist accounts of causation~\citep{pearl2009causality}:

\begin{definition}[Causal Observation Utility]
\label{def:causal}
The observations $\mathcal{O}(\tau)$ in a trajectory $\tau$ have \emph{causal utility} if, when they are removed from the reasoning context, the model fails to produce the correct answer:
\begin{equation}
    P_{\mathcal{M}}(\hat{a} = a^* \mid q, I, \tau \setminus \mathcal{O}(\tau)) < P_{\mathcal{M}}(\hat{a} = a^* \mid q, I, \tau).
\end{equation}
\end{definition}

In practice, we implement this test by presenting the student model $\mathcal{M}_S$ with the trajectory in which all tool observations are masked out (replaced by a placeholder), and checking whether it can still reach $a^*$.
If it can, the observations were dispensable and the trajectory is spurious; if it cannot, the observations were causally necessary.

\subsection{From Definition to Method}

The two conditions in Definition~\ref{def:learnable} directly motivate our \textbf{dual rejection sampling} procedure:
\begin{itemize}[leftmargin=*, itemsep=2pt]
    \item \textbf{Stage~1 (Outcome Rejection):} discard trajectories where $\hat{a} \neq a^*$ \;$\longrightarrow$\; enforces Condition~1.
    \item \textbf{Stage~2 (Observation Rejection):} among correct trajectories, discard those where the model still answers correctly without observations \;$\longrightarrow$\; enforces Condition~2.
\end{itemize}

The retained set consists exclusively of learnable trajectories—those that are both answer-correct and tool-beneficial.
Section~\ref{sec:method} details the full pipeline, including how trajectories are generated, how the counterfactual test is efficiently implemented, and how the framework scales across diverse visual reasoning domains.